\section{生成系统栈：DSL、Context、Agent/Workflow 与 System1}

前面讨论了为什么做视频，本章进入王冠提出的核心方法论：生成系统。生成系统不是单个模型，也不是单个 workflow，而是一套围绕某类信息商品的生产方法。它把问题表达、有效 context、agent/workflow、模型能力和评测组织成一个系统。

\begin{figure}[H]
\centering
\includegraphics[width=0.92\textwidth]{generation-system-stack.png}
\caption{生成系统栈：System1、DSL、Context、Agent/Workflow 和评测协同。自制概念图，依据 02:01:44--02:25:49 对谈内容整理。}
\end{figure}

\begin{knowledgebox}{读图：生成系统主要是 System2}
System1 是基座模型；DSL 用来表达问题和生产方法；Context 提供有效 token；Agent/Workflow 组织步骤；评测判断系统智慧程度。王冠认为 agent 和 workflow 不必对立，它们本质上都在产生更有效的 context，让 System1 输出更好。
\end{knowledgebox}

\subsection{DSL：把生产方法表达出来}

DSL 是 Domain-Specific Language，领域特定语言。这里不一定指传统编程语言，而是某个生成任务的表达体系。对视频来说，DSL 要能描述镜头、节奏、画面、风格、素材、转场、目标受众和生产方法。没有 DSL，系统只能堆 prompt；有了 DSL，系统才能把生产方法稳定复用。

\subsection{Context Engineering：产品经理的新主战场}

本节从 DSL 转向 context，因为生成系统的效果往往不只取决于模型本身，还取决于模型在生成前看到了什么、如何理解任务结构。

王冠把很多 agent、workflow、知识库和专库都看成 context 的生成机制。模型输出取决于前文 token，产品要做的是把业务 know-how 转成模型可理解、经济上划算、效果上稳定的 context。这个系统不一定技术上最复杂，却高度依赖业务理解和产品品味。

\begin{warningbox}{不要把 Workflow 低估成“没技术含量”}
Workflow、Agent、专库和 DSL 都是在给模型提供有效 context。它们的价值不只在工程复杂度，而在能否把业务方法、任务结构和评测目标稳定表达出来。
\end{warningbox}

\subsection{北极星指标：系统智慧程度}

本节回答“AI 产品到底优化什么”。商业价值是底线，但如果只看短期收入，就会错过生成系统作为生产能力本身的成长。

王冠认为 AI Native 产品当然要有商业价值，但更高层的北极星指标是“系统智慧程度”。这不是玄学，而是问：系统是否能自动生产、自动选择方法、自动修正、自动形成更好的输出？对生成系统来说，用户不只买一个工具，而是在使用一个会越来越会做事的生产系统。

\begin{figure}[H]
\centering
\includegraphics[width=0.96\textwidth]{ai-product-north-star.png}
\caption{AI 产品北极星：商业价值之上，用系统智慧程度衡量产品。自制概念图，依据 02:50:34--03:05:45 对谈内容整理。}
\end{figure}

\begin{knowledgebox}{读图：智慧程度不是抽象口号}
图中从商业价值到能力设计，再到 context、自动生产和智慧程度。系统智慧程度可以体现在：是否理解目标，是否选择合适方法，是否减少人工步骤，是否能从失败中改进，是否让用户以更低成本完成更高质量生产。
\end{knowledgebox}

\subsection{本章小结}

生成系统是应用公司可能建立壁垒的核心。它把模型能力转成生产方法，把业务 know-how 转成 context，把单次生成转成可复用 recipe，并用系统智慧程度衡量产品进步。

